\def\topic{Linear algebra}
\ifdefined\MyHeaderIncluded
\else
\def\qwertyflag{true}

\ifdefined\MyHeaderIncluded

\else
  \newcommand{\foobar}{FooBar}
  \def\MyHeaderIncluded{true}
% if is till the end of the file

\documentclass[10pt,a4paper]{article}
\usepackage[utf8]{inputenc}
%\usepackage[german]{babel}
%\usepackage[cyrillic]{babel}
\usepackage[T1]{fontenc}
%\usepackage[T2A]{fontenc}
%\usepackage[english, russian]{babel}
\usepackage{amsmath}
\usepackage{amsfonts}
\usepackage{amssymb}
\usepackage{amsthm}
\usepackage{mathtools} %coloneqq

\usepackage{bussproofs} %

%\renewcommand{\thesubsection}{\arabic{subsection}}
%\newtheorem{thm}{Theorem}[section]
%\renewcommand{\thethm}{\arabic{subsection}.\arabic{thm}}

\usepackage{thmtools}
\usepackage{hyperref}
\hypersetup{
    colorlinks=true,
    linkcolor=blue,
    filecolor=magenta,      
    urlcolor=cyan,
    pdftitle={Overleaf Example},
    %pdfpagemode=FullScreen,
    }
%\usepackage{stix-type1}%\turnediota
\theoremstyle{definition}

\newtheorem{constructor}{Constructor}[section]
\newtheorem{iaxiom}{Axiom}
\newtheorem{axiom}{Axiom}
\newtheorem*{axiom*}{Axiom}
\newtheorem{definition}{Definition}[section]
\newtheorem*{definition*}{Definition}
\newtheorem{remark}{Remark}[section]
\newtheorem*{remark*}{Remark}
\newtheorem{recall}{Recall}[section]
\newtheorem*{recall*}{Recall}

\newtheorem{lemma}{Lemma}[section]
\newtheorem{theorem}{Theorem}[section]
\newtheorem*{theorem*}{Theorem}
\newtheorem{note}{Note}[section]
\usepackage{etoolbox}

%\newcommand*{\@Dashv}[2]{%
%  \reflectbox{$\m@th#1#2$}%
%}
\newcommand*{\separation}[3]{% 1)binder, 2)set, 3)condition
  \{#1\in #2: #3\}
}

% https://www.freecodecamp.org/news/simple-regex-tricks-for-beginners-3acb3fa257cb/
%  \\mathcal\{T\} |---> \topology
%  \\mathcal\{([^\}]*)\}   |--->  $1
\newcommand{\AND}{\ \ \&\ \ }
\newcommand{\any}{\mathrm{\ any}}
\newcommand{\borel}{\mathcal{B}}
\newcommand{\can}{\mathrm{can}}
\newcommand{\canonize}{\mathrm{canonize}}
\newcommand{\closure}{\mbox{cl}}
\newcommand{\complex}{\mathbb{C}}
\newcommand{\converges}{\longrightarrow}
\newcommand{\countable}{\mbox{\ countable}}
\newcommand{\ldefined}{\mbox{\ defined}}
\newcommand{\dom}{\mbox{dom}}
\newcommand{\Dom}{\mbox{Dom}}
\newcommand{\Done}{\noindent \textbf{Done.}\\}
\newcommand{\lexists}{\mbox{\ exists}}
\newcommand{\lforall}{\mbox{\ forall}}
\newcommand{\lunique}{\mbox{\ unique}}
\newcommand{\family}{\mathrm{\ family}}
\newcommand{\Fm}{\mathrm{Fm}}
\newcommand{\FuncSymb}{\mathrm{FuncSymb}}
\newcommand{\function}{\mathrm{\ function}}
\newcommand{\functional}{\mbox{\ functional}}
\newcommand{\IFF}{\ \Longleftrightarrow\ }
\renewcommand{\iff}{\ \longleftrightarrow\ }
\newcommand{\IMPLIES}{\ \Longrightarrow\ }
\renewcommand{\implies}{\ \longrightarrow\ }
\newcommand{\Inf}{\mathrm{Inf}}
\newcommand{\injective}{\mbox{\ injective}}
\newcommand{\injection}{\mbox{\ injection}}
\newcommand{\interior}{\mbox{int}}
\newcommand{\intersection}{\bigcap}
\newcommand{\isempty}{\mathrm{\ empty}}
\newcommand{\LB}{\mathrm{LB}}
\newcommand{\nghbset}{\mathcal{N}}
\newcommand{\nonempty}{\mathrm{\ nonempty}}
\newcommand{\Ord}{\mathrm{Ord}}
\newcommand{\power}{\mathcal{P}}
\newcommand{\pr}{\mathrm{pr}}
\newcommand{\fst}{\mathrm{fst}}
\newcommand{\snd}{\mathrm{snd}}
\newcommand{\PredSymb}{\mathrm{PredSymb}}
\newcommand{\Proof}{\noindent \textbf{Proof.}\\}
\newcommand{\ran}{\mathrm{ran}}
\newcommand{\Ran}{\mathrm{Ran}}
\newcommand{\real}{\mathbb{R}}
\newcommand{\relation}{\mathrm{\ relation}}
\newcommand{\set}{\mathrm{\ set}}
\newcommand{\singleton}{\mathrm{\ singleton}}
\newcommand{\semantics}{\mathrm{semantics}}
\newcommand{\semialgebra}{\mathcal{S}}
\newcommand{\Sup}{\mathrm{Sup}}
\newcommand{\surjective}{\mbox{\ surjective\ on\ }}
\newcommand{\Term}{\mathrm{Term}}
\newcommand{\Type}{\mathrm{Type}}
\newcommand{\topology}{\mathcal{T}}
\newcommand{\total}{\mbox{\ total\ on\ }}
\newcommand{\transitive}{\mathrm{\ transitive}}
\newcommand{\union}{\bigcup}
\newcommand{\lTrue}{\mathrm{True}}
\newcommand{\lFalse}{\mathrm{False}}
\newcommand{\lUncertain}{\mathrm{Uncertain}}
\newcommand{\lif}{\mathrm{if}}
\newcommand{\lthen}{\mathrm{then}}
\newcommand{\lelse}{\mathrm{else}}
\newcommand{\lexception}{\mathrm{exception}}
\newcommand{\lundef}{\mathrm{undef}}
\newcommand{\univocal}{\mathrm{\ univocal}}
\newcommand{\univocalat}{\mathrm{\ univocal\ at\ }}
\newcommand{\UB}{\mathrm{UB}}
\newcommand{\Un}{\mathcal{U}}
\newcommand{\Var}{\mathrm{Var}}

\newcommand{\notvarphi}{\phi}
\renewcommand{\phi}{\varphi}

\newcommand{\comment}[1]{\quad\mbox{(#1)}}
% local commands
%\preto\section{\undef\blipo} %etoolbox
\preto\subsection{\undef\blipo} %etoolbox
% https://tex.stackexchange.com/questions/105684/limit-macro-newcommand-scope-to-section
\usepackage{graphicx}
\usepackage{authblk}

\setlength\parindent{0pt} % whole doc of \noindent
\usepackage{soul}
%\usepackage{mathgl}

\fi % guard
\ifdefined\topic
\else
\def\topic{No Titel}
\fi
\title{\topic}
\author[1]{Georgy Dunaev}
\begin{document}
\maketitle
\fi
\section{Linear algebra}
These preliminary definitions describe matrices and the combination of
maps before imposing algebraic operations on their values.

\begin{definition}[Matrix over a set]
A matrix $A$ over a set $S$ is a function
\[
A:\{1,\ldots,m\}\times\{1,\ldots,n\}\longrightarrow S
\]
for specified nonnegative integer dimensions $m,n$. Its entry in row $i$ and column $j$ is
$A(i,j)$. If zero dimensions are allowed, the corresponding index set is
empty. Addition and multiplication of matrices require further structure
on $S$.
\end{definition}

\begin{definition}[Combining matching graph entries]
For functions $f_1$ and $f_2$, define
\[
\mbox{combine}(f_1,f_2)
 \ceq\cc{(((m,k),(m,v)),(m,(k,v)))}
 {(m,k)\in f_1\land(m,v)\in f_2}.
\]
\end{definition}
This relation combines two graph entries with the same input $m$ into a
single entry with paired output $(k,v)$. When $f_1:M\to K$ and
$f_2:M\to L$, it records the graph-level construction of the map
$m\mapsto(f_1(m),f_2(m))$ into $K\times L$. This is the pairing of maps
with a common domain in the category of sets and functions.

\subsection{Finite-dimensional normed spaces}
The following results use set-sized real or complex vector spaces in the
classical analytic setting. Completeness of $\mathbb R$ and $\mathbb C$
is assumed. The metric compactness criterion from the topology chapter
supplies the only compactness implication used below. Its classical
countable-choice convention is retained here; in particular, the proof
of coordinate-sphere compactness uses that criterion.

\begin{lemma}[Compact coordinate spheres]
For $n\geq1$, the set
\[
S_n=\{a\in\mathbb K^n:\max_j|a_j|=1\},
\qquad \mathbb K\in\{\mathbb R,\mathbb C\},
\]
is compact in the coordinate maximum metric.
\end{lemma}
\Proof
A Cauchy sequence of vectors is Cauchy in each coordinate; scalar
completeness gives a coordinatewise limit, which is also a limit in the
maximum metric because there are only finitely many coordinates.
The maximum norm is continuous, so $S_n$ is closed and complete.
Every bounded coordinate box is totally bounded: partition a containing
real interval into finitely many intervals of length less than a prescribed
$\varepsilon$, and take the finite product grid. In the complex case do
this for real and imaginary parts, with mesh small enough to bound the
complex modulus. A subset inherits total boundedness, by keeping only
balls that meet the subset and recentering them at such points with
twice the radius. Thus $S_n$ is totally bounded. The metric compactness
criterion gives compactness.
\Done

\begin{theorem}[Coordinate estimates and equivalence of norms]
\label{thm:finite_dimensional_norms}
Let $e_1,\ldots,e_n$ be a basis of a nonzero finite-dimensional normed
space $E$ over $\mathbb K$. There are constants $c,C>0$ such that
\[
c\max_j|a_j|\leq\Bigl\|\sum_{j=1}^n a_je_j\Bigr\|
\leq C\max_j|a_j|\qquad(a\in\mathbb K^n).
\]
Consequently any two norms on a finite-dimensional vector space are
equivalent and give the same topology.
\end{theorem}
\Proof
For the upper bound take $C=\sum_j\|e_j\|>0$ and use the triangle
inequality. The map $a\mapsto\|\sum_j a_je_j\|$ is continuous in the
coordinate maximum metric, since the reverse triangle inequality bounds
its change by $C\max_j|a_j-b_j|$. On $S_n$ it is strictly positive,
by linear independence. Compactness supplies a positive lower bound:
each point has a neighbourhood on which the function exceeds half its
value there; finitely many such neighbourhoods cover $S_n$, and the
minimum of their finitely many positive bounds is positive. Call it $c$.
For $a\ne0$, apply this bound to $a/\max_j|a_j|$ and rescale. The
zero vector satisfies both inequalities directly.

Apply these estimates separately to two norms, using the same basis;
each norm bounds and is bounded by positive multiples of the coordinate
maximum norm, hence of the other norm. The zero-dimensional space has
only its zero vector, so the conclusion there is immediate.
\Done

\begin{corollary}[Completeness and closedness in finite dimensions]
Every finite-dimensional normed space is complete. Every finite-dimensional
linear subspace of an arbitrary normed space is closed.
\end{corollary}
\Proof
For a Cauchy sequence in a finite-dimensional space, the lower coordinate
estimate makes each coordinate sequence Cauchy. Scalar completeness gives
limits for all coordinates, and the upper estimate proves convergence to
their linear combination. This also covers the zero-dimensional space.
Now let $M$ be a finite-dimensional subspace of a normed space $E$,
and take $x\notin M$. Adjoin $x$ to a basis of $M$ and apply the lower
coordinate estimate in the finite-dimensional span of $M\cup\{x\}$.
The coefficient of $x$ in $x-m$ is one for every $m\in M$, so
$\|x-m\|\geq c>0$ for a fixed $c$. Thus the ball of radius $c$ about
$x$ misses $M$. Every point outside $M$ has such a neighbourhood,
which proves that $M$ is closed without a successive-choice argument.
\Done

The weak topology is the coarsest topology making every continuous
linear functional continuous. Basic neighbourhoods impose finitely many
scalar inequalities; the functional-analysis chapter develops this
definition further.

\begin{corollary}[Finite-dimensional compactness and weak topology]
In a finite-dimensional normed space, every closed bounded set is compact,
and the weak topology agrees with the norm topology.
\end{corollary}
\Proof
The coordinate estimates carry bounded sets into bounded coordinate boxes
and carry finite coordinate nets back to finite norm nets. Hence a bounded
set is totally bounded. A closed set is complete by the preceding
corollary, so closed bounded sets are compact by the metric criterion.

Each coordinate functional $\ell_j(\sum_k a_ke_k)=a_j$ is continuous,
since $|\ell_j(x)|\leq\|x\|/c$. For every $\varepsilon>0$ the weakly
open set
\[
\{x:|\ell_j(x)|<\varepsilon/C\text{ for }1\leq j\leq n\}
\]
is contained in the norm ball of radius $\varepsilon$ by the upper
estimate. Translations give the same conclusion at any centre, so every
norm-open set is weakly open. Conversely every continuous linear
functional is norm-continuous by definition, so the weak topology is
coarser than the norm topology. They are equal. In the zero-dimensional
case there is only one topology.
\Done

\begin{corollary}[Linear maps and finite-dimensional range]
Every linear map from a finite-dimensional normed space $E$ into a
normed space $F$ is bounded. More generally, if a linear map $T:E\to F$
is bounded and its range is finite-dimensional, then the image of the
closed unit ball has compact closure in $F$. Neither space needs to be
complete in this latter assertion.
\end{corollary}
\Proof
If $E\ne\{0\}$, choose a basis $e_1,\ldots,e_n$ and the lower
coordinate constant $c>0$. Linearity and the triangle inequality give
\[
\left\|T\left(\sum_j a_je_j\right)\right\|
\leq \left(\sum_j\|Te_j\|\right)\max_j|a_j|
\leq \frac{\sum_j\|Te_j\|}{c}\left\|\sum_j a_je_j\right\|.
\]
This is boundedness; a map on the zero space is zero and is bounded too.
For the second assertion, the range $M=T[E]$ is a closed
finite-dimensional subspace of $F$. Boundedness of $T$ makes the unit-ball
image bounded in $M$. Its closure in $F$ stays in $M$ and remains bounded;
it is closed in $M$, hence compact by finite-dimensional compactness.
\Done

The second assertion explains why bounded finite-rank operators are
compact, where an operator is called compact when its unit-ball image
has compact closure. Finite rank controls the size of the range;
boundedness separately controls the images of bounded sets.

\subsection{The unit ball detects finite dimension}
\begin{lemma}[Riesz's separation lemma]
Let $M$ be a proper closed linear subspace of a normed space $E$, and
let $0<\theta<1$. There is a unit vector $x$ with
\[
\inf_{m\in M}\|x-m\|>\theta.
\]
\end{lemma}
\Proof
Choose $z\notin M$. Closedness gives
$d=\inf_{m\in M}\|z-m\|>0$. By the definition of infimum, choose
$m_0\in M$ with $\|z-m_0\|<d/\theta$, and put
$x=(z-m_0)/\|z-m_0\|$. For any $m\in M$,
\[
\|x-m\|=
\frac{\|z-(m_0+\|z-m_0\|m)\|}{\|z-m_0\|}
\geq\frac d{\|z-m_0\|}>\theta.
\]
The same lower bound holds for the infimum, proving the claim.
\Done

\begin{theorem}[Compact unit balls and dimension]
In the usual classical setting with countable dependent choices, the
closed unit ball of a normed space is compact if and only if the space
is finite-dimensional. No completeness assumption on the space is needed.
\end{theorem}
\Proof
The finite-dimensional implication is the preceding compactness corollary.
If $E$ is infinite-dimensional, start with any unit vector $x_1$.
Having chosen $x_1,\ldots,x_n$, their span is finite-dimensional,
closed, and proper. Riesz's lemma with $\theta=1/2$ supplies a unit
vector $x_{n+1}$ at distance greater than $1/2$ from that span.
Consequently $\|x_i-x_j\|>1/2$ for all distinct indices. A ball of
radius $1/4$ cannot contain two of these vectors, by the triangle
inequality. The closed unit ball therefore has no finite cover by
such balls and is not totally bounded, hence not compact.
The sequence of successive selections is the place where the stated
choice convention is used; finite-dimensional closedness did not
require this infinite construction.
\Done

This theorem concerns the norm topology. Infinite-dimensional closed
unit balls may be compact in a weaker topology; the distinction is
central to the functional-analysis chapter. Nor does closedness of a
subspace alone imply finite dimension: the finite-dimensional hypothesis
above is a sufficient condition for closedness, not a characterization.

\ifdefined\qwertyflag
\undef\qwertyflag
\end{document}
\fi
